\subsection{Write the selected company rules and commands in one source}
\label{sec:selected-company-source}

A generated method can calculate the right journal while leaving the modeller
unable to state the rule that generated it. In the claim example, the customer
check, the source of a claim identifier and the order of assessment updates are
substantive choices. Naming a template by its digest records which choice was
made, but does not let an author write or change that choice. The selected
company source therefore includes the opening conditions, state invariant,
eligibility stages and construction recipes themselves. The earlier source
frontends remain useful predecessor comparisons; this successor targets the
declared opening and execution path together. The opening template in
this profile still requires an empty claim map. Section~\ref{sec:complete-claims-opening}
extends that boundary for explicitly complete supplied carried claims.

Consider a collection whose claim identifier is absent. Reading that claim's
customer immediately would fail for a missing row, obscuring the intended
refusal that the claim does not exist. In \code{OrderedEligibilityText/1}, the
author writes a membership stage before the customer stage. A binding such as
\code{(map\_contains "/claims" "claim\_id")} supplies the membership
predicate. A later \code{map\_field} binding names the same map and event
input, the customer field and the preceding presence stage. Compilation accepts
that lookup only when the earlier stage unconditionally requires the matching
membership fact. All stages compile, including ones that an earlier refusal
would prevent from running. An empty map cannot hide a malformed field lookup.

The full event-input types, stage identifiers, failure reasons and guarded
predicates are also declarations. The retained order first resolves the outer
event structure and unknown inputs, then approval and the operation's specific
conditions. Moving claim existence ahead of the sealed-day condition changes
which reason is returned when both fail. Negating the approval predicate
changes which events are eligible. These edits demonstrate that the source
controls the selected rule; they are different rules, not equivalent ways to
spell the retained contract. The six eligibility sources reconstruct their
complete typed templates and preserve 38 eligible and 34 refused gate results.

Opening admission has a related but distinct job. \code{OpeningAdmissionText/1}
states structural checks for the input fields, policy values, row and clock
schemas, roles and report mapping, alongside typed predicates for the dates
and amounts. The following excerpt is an actual declaration for the final
command-count condition; its position follows the other opening checks.

\begin{quote}\small\ttfamily
stage "commands" commands "command\_limit";\\
check 64;\\
end\_stage;
\end{quote}

Other structural stages declare their paths and expected values directly.
Predicate stages use the same guarded-expression grammar as eligibility.
Dependency declarations name the exact reused semantic implementations, while
identities for locally authored predicates and templates are computed from the
definitions. The two 21-stage sources preserve all 82 complete opening results.
Ten inputs are admitted, including those whose future command will later fail.
Neither successful parsing nor admitted opening conditions establish that the
whole future program succeeds.

Construction recipes express the next step: how a permitted event supplies the
values needed by a method. \code{MethodRecipeText/1} keeps the existing
method declarations for monetary calculations, postings, cash flows, guarded
conditions and memory writes, and adds explicitly typed parameters. For
example, a customer parameter can read the event's customer input, and a
memory write can refer to that parameter. No event value is interpreted as
source code. A pointer parameter combines literal path segments and previously
declared text parameters; the existing pointer encoder escapes slashes and
tildes in a claim identifier before using it as a path.

The distinction matters for the sale and collection methods. The sale writes
the customer's name, due date, event identifier and issue date into a new
claim. The collection reads the selected claim's remaining amount through a
parameterized pointer. These values are supplied by the actual event and
current book. The accounting expressions and account choices remain visible
in the recipe. An author need not reproduce a serialized method tree or
maintain its selection hash manually. Missing, unused, forward or mistyped
parameters are refused by the original contextual constructor.

An assessment requires a sequence of writes over the current claims.
\code{OrderedExpansionText/1} contains the complete base recipe, map path,
row types, per-row replacements and final replacements. Each replacement
declares its identifier, path and value expression. A row expression can read
a field, the current key, an event input or event date, and choose between
typed branches of a date condition. Both branches are checked even when only
one will be evaluated. The explicit assessment-key order determines the write
sequence; a duplicate or incomplete order refuses construction. A malformed
row expression also refuses when the map happens to be empty.

Four scalar recipe sources and two ordered-expansion sources reconstruct their
complete templates. Across 18 scalar and 20 expansion cases, the full
construction results, transition requests and transitions equal the retained
typed comparisons. These checks include identifiers containing slashes, tildes
and Unicode. Source lowering reports syntax and structure only; the current
book, event, schemas and memory definition remain necessary for contextual
construction. Thus a reusable recipe is not prematurely described as a method
valid for every possible state.

\code{SelectedCompanyText/1} joins these declarations. It imports the opening
book, schemas, statement mapping and memory definition by exact content
identity. Those four resources contain facts and types. The opening rule,
invariant, eligibility and construction algorithms are inline source blocks.
Roles, unit-cost pointer and policy define the selected company conventions.
Events retain their full typed inputs, provenance and declaration order;
reports and carry commands retain their dates, notes, equity movements and
explicit memory policies. The native-shaped admission input is derived from
these same facts and projected commands, rather than authored as an unrelated
second copy.

Lowering creates the full program and its selection without running commands.
Execution then reuses the opening-to-program binding and the existing
eligibility, invariant, construction, transition, report and carry evaluators.
For four successful programs, all 44 complete command results are preserved.
For 26 adverse programs, the stopped prefix, current and opening states,
refused request where present and unexecuted suffix are preserved as well.
All native lifecycle functions and the operating-method builder are disabled
during these comparisons. Changing the first sale's approval directly in the
company source stops the program before any command succeeds.

One provenance distinction must remain visible. Canonical content identities
ignore dictionary insertion order, whereas the outer execution binding also
records the serialized order of the complete supplied program. The source
parsers reconstruct some definition dictionaries in a different order. Their
canonical content is identical, but that one outer order digest changes. The
comparison recomputes and retains both order identities, and permits a
difference only at \code{program\_order\_sha256}. Every other result field
matches, including the separate digest of event-input name order, which must
remain unchanged because it affects the order of reported unknown inputs.
This is a precise provenance exception, not a relaxation of financial-result
or failure comparison.

Source identity retains the original UTF-8 bytes. Whitespace and CRLF line
endings can change that identity while leaving the lowered program and its
results equal. Errors inside a nested rule or recipe identify the corresponding
bytes in the full company source. A lexical error can identify its token; a
semantic error that is known only at stage or document scope is reported at
that scope. The corrected nested-invariant reader described in
Section~\ref{sec:complete-claims-opening} also preserves quoted Unicode
separator characters as data rather than treating them as source line endings. The checked sources reject duplicate or missing declarations,
changed or unused imported resources and incompatible operation selections.

The comparison supports this finite synthetic company language. It does not
make an edited rule accounting-equivalent to the retained rule, supply missing
facts about a real company's claim population, or complete the broader fiscal,
sector and measurement algorithms. Independent implementability and reader
acceptance remain open. The complete source examples and result pairs are
retained under
\begingroup\raggedright
\path{data/accounting-master/selected-company-source-20261026/}.
The component evidence is under
\path{data/accounting-master/ordered-eligibility-source-20261026/},
\path{data/accounting-master/opening-admission-source-20261026/}
and \path{data/accounting-master/method-recipe-source-20261026/}.
\par\endgroup
